\section{classes of regularity}

Let $A\subset \RR$ and let $f\colon A\to \RR$ be a function defined on $A$.
We will say that $f$ is of class $C^0$ if $f$ is continuous.
We can define $C^0(A)$ as the set of continuous functions defined on $A$ (with
values in $\RR$).
If $f$ is differentiable (on all of $A$) and the derivative of $f$ is also
continuous (on all of $A$), we will say that $f$ is of class $C^1$ and we will
write $f\in C^1(A)$.
Continuing: if $f$ is differentiable and $f'$ is also differentiable and if the
second derivative $f''$ is continuous, we will say that $f$ is of class $C^2$.
And so on.

It is important to note that $C^1(A)$ does not coincide with the set of
differentiable functions on $A$.
Indeed, we have already seen in example~\ref{ex:derivata_non_continua} that
there exist differentiable functions whose derivative is not continuous and
therefore such functions, although differentiable, are not of class $C^1$.

To formally define the set of functions of class $C^k$ for each $k\in \NN$, we
must provide a definition by induction.
We also define the class $C^{\infty}(A)$, which is the set of functions that are
differentiable $k$ times for every $k\in \NN$.

\begin{definition}[functions of class $C^k$]
\mymark{**}
Given $A\subset \RR$, we denote by $\RR^A$ the set of all functions $f\colon A
\to \RR$.
For each $k\in \NN$, we define $C^k(A) \subset \RR^A$ as follows:
\mymargin{$C^k$}%
\index{$C^k$}
\begin{enumerate}
\item if $k=0$, we set
\[
  C^0(A) = \ENCLOSE{f\in \RR^A \colon \text{$f$ continuous}}.
\]
\item
  if $k>0$, we define inductively
  \[
    C^{k}(A) = \ENCLOSE{f\in \RR^A \colon \text{$f$ differentiable and $f'\in
  C^{k-1}(A)$}}.
  \]
\end{enumerate}
Clearly, if $j\ge k$, we have $C^j(A) \subset C^k(A)$, thus $C^k(A)$ is a
decreasing family (with respect to set inclusion), and it makes sense to define:
\mymargin{$C^\infty$}%
\index{$C^\infty$}
\[
    C^\infty(A) = \ENCLOSE{f\in \RR^A\colon \forall k \in \NN\colon f\in C^k(A)}
  = \bigcap_{k\in \NN} C^k(A).
\]

Functions $f\in C^k(A)$ are differentiable $k$ times.
We use the notation $f^{(j)}$%
\mymargin{$f^{(j)}$}%
\index{$f^{(j)}$} to denote the $j$-th derivative of a function $f$.
Thus, we have
\[
  f^{(0)} = f, \qquad
  f^{(1)} = f', \qquad
  f^{(2)} = f'', \dots
\]
\end{definition}

We have already observed that $\RR^A$ is a real vector space.
The spaces $C^k(A)$ for $k=0, \dots, \infty$ are a decreasing family of vector
subspaces of $\RR^A$.
Indeed, we know that the linear combination of continuous functions is
continuous and that the linear combination of differentiable functions is
differentiable.

\begin{theorem}
The functions $x$, $e^x$, $\ln x$, $\sin x$, $\cos x$, $\tan x$, $\arctan x$ are
of class $C^\infty$.
Sum, difference, product, quotient, and composition of functions of class $C^n$
are also of class $C^n$ for $n=0,1,2,\dots,\infty$.
\end{theorem}
\begin{proof}
We verify by induction that if $f,g \in C^n$, then $f+g \in C^n$.
For $n=0$, we know that the sum of continuous functions is continuous: $f,g\in
C^0$ implies $f+g\in C^0$.
For $n>1$, if $f,g\in C^n$, then $f$ and $g$ are differentiable, and we have
$(f+g)'=f'+g'$ (theorem~\ref{th:derivate_operazioni}).
But by definition of $C^n$, we have $f',g'\in C^{n-1}$, and by the inductive
hypothesis, $f'+g'\in C^{n-1}$, thus $f+g\in C^n$.

A similar proof can be made for the product $f\cdot g$ by observing that
$(f\cdot g)' = f'g + fg'$ and that if $f,g\in C^n$, then $f',g'\in C^{n-1}$,
thus $f'g,fg'\in C^{n-1}$ and $f'g+fg'\in C^{n-1}$.

The same holds for composition since $(f\circ g)' = (f'\circ g)\cdot g'$.

Thus, the sum, product, and composition of functions of class $C^n$ are of class
$C^n$, and the same holds for functions of class $C^\infty$.

Now note that the functions $-x$ and $\frac 1 x$ are of class $C^\infty$ (we can
explicitly compute their derivatives of any order), and therefore if $f\in C^n$,
then $-f$ and $\frac 1 f$ are of class $C^n$ due to the property of composition.
It follows that the difference and the quotient of functions of class $C^n$ are
also of class $C^n$.

We can then consider the elementary functions listed in the statement and
observe that they are all differentiable functions and that their derivative can
be expressed as a sum, product, quotient, or composition of the same elementary
functions.
Thus, if all these functions are of class $C^n$, then their derivative is also
of class $C^n$.
But if the derivative is of class $C^n$, then the function is of class
$C^{n+1}$, and thus, by induction, all elementary functions are of class
$C^\infty$.
\end{proof}

\begin{example}
  The function 
  \[
    \frac{e^{x-\arctan x^2}}{1+\ln x}
  \]
  is of class $C^\infty$.
\end{example}

\begin{definition}[Lipschitz functions]
\mymark{**}
\mymargin{Lipschitz function}
\index{function!Lipschitz}
\index{Lipschitz}
A function $f\colon A\subset \RR \to \RR$ is said to be \emph{Lipschitz} (or
$L$-Lipschitz if we want to emphasize the dependence on $L$) if there exists
$L>0$ such that for every $x,y\in A$, we have
\[
  \abs{f(x) - f(y)} \le L\cdot  \abs{x-y}.
\]
The smallest constant $L$ for which the previous relation is satisfied is called
the \emph{Lipschitz constant}
\mymargin{Lipschitz constant}
\index{constant!Lipschitz}
of $f$.
\end{definition}

\begin{theorem}[Lipschitz criterion]
\mymark{**}
Let $A\subset \RR$ and let $f\colon A \to \RR$ be a Lipschitz function with
Lipschitz constant $L$.
If $f$ is differentiable at a point $x\in A$, then $\abs{f'(x)}\le L$.
Conversely, if $f\colon I \subset \RR \to \RR$ is a differentiable function
defined on an interval $I$ and if there exists $L$ such that for every $x\in I$,
we have $\abs{f'(x)}\le L$, then $f$ is Lipschitz.
\end{theorem}
%
\begin{proof}
If $f$ is Lipschitz, it means that the difference quotient is bounded.
That is, there exists $L>0$ such that
\[
  \abs{\frac{f(x) - f(y)}{x-y}} \le L \qquad \forall x,y\in A.
\]
Thus, the derivative, which is the limit of the difference quotient, if it
exists, is also bounded by the same constant: $\abs{f'(x)}\le L$ for every $x
\in A$.

Conversely, if the derivative is bounded $\abs{f'(z)}\le L$ for every $z \in I$
and if $x,y\in I$ are any points, then, by the Mean Value Theorem, the
difference quotient of $f$ is equal to the derivative at some point $z\in(x,y)$:
\[
  \abs{\frac{f(x) - f(y)}{x-y}} = \abs{f'(z)} \le L.
\]
and thus the function is $L$-Lipschitz:
\[
  \abs{f(x)- f(y)} \le L \abs{x-y}.
\]
\end{proof}

\begin{definition}[Hölder functions]
Let $\alpha>0$.
\mymargin{Hölder function}
\index{function!Hölder}
\index{Hölder}
A function $f\colon A \subset \RR \to \RR$ is said to be $\alpha$-Hölder if
there exists a constant $C>0$ such that
\begin{equation}\label{eq:2964536}
  \abs{f(x) - f(y)} \le C \abs{x-y}^\alpha.
\end{equation}
\end{definition}

\begin{definition}[uniform continuity]
\mymark{**}%
A function $f\colon A \subset \RR \to \RR$ is said to be \emph{uniformly
continuous}
\index{uniform continuity}%
\index{continuity!uniform}%
\mymargin{uniform continuity}%
if
\[
 \forall \eps>0\colon \exists \delta > 0 \colon
 \forall x,y\in A \colon \abs{x-y} < \delta \implies \abs{f(x)-f(y)} < \eps.
\]
\end{definition}

\begin{theorem}
Every Lipschitz function is $1$-Hölder (and vice versa).
Every $\alpha$-Hölder function is uniformly continuous.
Every uniformly continuous function is continuous.
Every $\alpha$-Hölder function with $\alpha>1$ has a zero derivative.
\end{theorem}
%
\begin{proof}
The first three statements follow directly from the definitions.
For the last observation, note that if $\alpha>1$ in
inequality~\eqref{eq:2964536}, one can divide by $\abs{x-y}$ and thus the
difference quotient tends to zero as $y\to x$.
\end{proof}

\begin{definition}[modulus of continuity]
Let $f\colon A\subset \RR \to\RR$ be a function.
We define the \emph{modulus of continuity}%
\mymargin{modulus of continuity}%
\index{modulus!of continuity} of $f$ as the function $M\colon$ $[0,+\infty) \to [0,+\infty]$ defined by
\[
  M(r) = \sup \ENCLOSE{\abs{f(x)-f(y)}\colon x,y \in A, \abs{x-y} \le r}.
\]
\end{definition}

\begin{theorem}[properties of the modulus of continuity]
Let $f\colon A\to \RR$ and let $M\colon [0,+\infty)\to [0,+\infty)$ be its
modulus of continuity.
The function $M(r)$ is increasing.

The function $f$ is uniformly continuous if and only if
\[
  \lim_{r\to 0} M(r) = 0.
\]

The function $f$ is Lipschitz if and only if there exists $L$ such that
\[
  M(r) \le Lr.
\]

The function $f$ is $\alpha$-Hölder if and only if there exists $C$ such that
\[
  M(r) \le C r^\alpha.
\]
\end{theorem}
%
\begin{proof}
Observe that the condition
\[
   M(r) \le c
\]
is equivalent to
\[
\forall x,y\in A \colon \abs{x-y}\le r \implies  \abs{f(x)-f(y)} \le c.
\]
The condition $M(r)\to 0$ as $r \to 0$ means
\[
 \forall \eps>0\colon \exists \delta>0 \colon \forall r>0\colon r<\delta \implies M(r)<\eps
\]
and thus becomes the condition of uniform continuity.

The conditions of Lipschitz continuity and $\alpha$-Hölder continuity also
follow immediately.
\end{proof}

\begin{theorem}[restriction / gluing of uniformly continuous functions]
Let $f\colon A \subset \RR \to \RR$ be a uniformly continuous function. If
$B\subset A$, the restriction $f_{|B}$ of $f$ to $B$ is also uniformly
continuous.

Let $I,J\subset \RR$ be intervals such that $I\cap J \neq \emptyset$.
Let $f\colon I \cup J \to \RR$ be a function. If $f_{|I}$ and $f_{|J}$ are
uniformly continuous, then $f$ is uniformly continuous.
\end{theorem}
\begin{proof}
The first part, regarding the restriction of a uniformly continuous function,
follows directly from the definition: if any property holds for every $x,y\in
A$, then it certainly holds for every $x,y\in B$ when $B\subset A$.

For the second part of the statement: suppose $f$ is uniformly continuous on $I$
and on $J$.
Let $\eps>0$ be given, and let $\delta_1$ and $\delta_2$ be the values of
$\delta$ given by the conditions of uniform continuity of $f$ on $I$ and on $J$,
respectively. Consider $\delta = \min\ENCLOSE{\delta_1, \delta_2}$.
Now let $x,y$ be any points in $I\cup J$ with $\abs{x-y}< \delta$.

There are two possible cases:
$x$ and $y$ are in the same interval ($I$ or $J$) or they are in different
intervals, one in $I$ and the other in $J$.
In the first case, since $\delta < \delta_1$ and $\delta< \delta_2$, the uniform
continuity of $f$ on $I$ and on $J$ ensures that $\abs{f(x)-f(y)} < \eps$ in any
case. In the second case, there must exist a point $z \in I\cap J$ that is an
intermediate point between $x$ and $y$.
Then, using the triangle inequality:
\[
  \abs{f(x) - f(y)} \le \abs{f(x) - f(z)} + \abs{f(z) - f(y)}
     \le \eps + \eps.
\]
thus (except for replacing $\eps$ with $\eps/2$) in this case as well, the
desired estimate is obtained.
\end{proof}

\begin{theorem}[Heine-Cantor]
  \label{th:heine_cantor}%
  \mymark{***}%
  \index{theorem!Heine-Cantor}%
  \mymargin{Heine-Cantor}%
Let $a,b\in \RR$, $a<b$.
Let $f\colon [a,b]\to \RR$ be a continuous function.
Then $f$ is uniformly continuous.
\end{theorem}
%
\begin{proof}
\mymark{***}%
Suppose for contradiction that $f$ is not uniformly continuous. Then $f$
satisfies the negation of the property
\[
 \forall \eps>0\colon \exists \delta > 0 \colon
 \forall x,y\in [a,b] \colon \abs{x-y} < \delta \implies \abs{f(x)-f(y)} < \eps
\]
which is
\[
 \exists \eps>0\colon \forall \delta > 0 \colon
 \exists x,y \in [a,b] \colon \abs{x-y} < \delta \land \abs{f(x)-f(y)} \ge \eps.
\]
Thus, given $\eps>0$ that satisfies the previous property, we can take
$\delta=1/k$ for each $k\in \NN$, obtaining two sequences $x_k$, $y_k$ such that
\[
  \abs{x_k-y_k} < \frac 1 k \qquad\text{and}\qquad \abs{f(x_k) - f(y_k)} \ge \eps.
\]
By the Bolzano-Weierstrass theorem, there will exist a convergent subsequence
$x_{k_j} \to c$.
Since $\abs{x_k-y_k}\to 0$, we must also have $y_{k_j}\to c$.
But then, by the continuity of $f$:
\[
 \abs{f(x_{k_j})-f(y_{k_j})} \to \abs{f(c) - f(c)} = 0
\]
in contradiction with the condition $\abs{f(x_k)-f(y_k)}\ge \eps$.
\end{proof}

\begin{theorem}[asymptote]
\index{theorem!asymptote}
\mymargin{asymptote theorem}
Let $f\colon [a,+\infty) \to \RR$ be a continuous function and let $g\colon
[a,+\infty) \to \RR$ be a uniformly continuous function such that
\[
  \lim_{x\to +\infty} f(x) - g(x) = 0.
\]
Then $f$ is also uniformly continuous.

In particular, if $f$ has an \emph{oblique asymptote}%
\mymargin{oblique asymptote}%
\index{asymptote!oblique} that is, if there exist $m\in \RR$ and $q\in \RR$ such that
\[
  \lim_{x\to +\infty} f(x) - (mx + q)  = 0
\]
then $f$, if continuous, is uniformly continuous.

A similar result holds for functions defined on intervals of the type
$(-\infty,a]$ (taking limits at $-\infty$) and thus for functions defined on all
of $\RR$ (taking limits both at $+\infty$ and $-\infty$).
\end{theorem}

\begin{proof}
For every $\eps>0$, there exists $M>a$ such that $\abs{f(x)-g(x)} < \eps/3$ for
every $x\ge M$. On the other hand, the function $f$, by the Heine-Cantor
theorem, is uniformly continuous on $[a,M+1]$, and thus there exists
$\delta_1>0$ such that for $x,y \in [a,M+1]$ with $\abs{x-y}< \delta_1$, we have
$\abs{f(x)-f(y)} < \eps$. On the other hand, $g$ is uniformly continuous on
$[M,+\infty)$, and therefore there exists $\delta_2$ such that for $x,y\in
[M,+\infty)$ with $\abs{x-y} < \delta_2$, we have $\abs{g(x)-g(y)} < \eps/3$.
But in this latter case, we have:
\begin{align*}
  \abs{f(x)-f(y)} &\le \abs{f(x) - g(x)} + \abs{g(x) - g(y)} + \abs{g(y)-f(y)} \\
  & \le \frac{\eps}{3} + \frac{\eps}{3} + \frac{\eps}{3} = \eps.
\end{align*}
Thus, setting $\delta = \min\ENCLOSE{1, \delta_1, \delta_2}$, for any $x,y\in
[a,+\infty)$ with $\abs{x-y}< \delta$, we are certainly in one of the two
previous cases, and therefore, in any case, we obtain $\abs{f(x)-f(y)} < \eps$,
as we needed to prove.

In the particular case $g(x) = mx +q$, it is simply observed that $g$ is
uniformly continuous since it is $L$-Lipschitz with $L=\abs{m}$.
\end{proof}